% ------------------------------------------------------------
% Elbische FDI-Schriften
% ------------------------------------------------------------
% Wird von \FDIInputLanguageFonts in der Praeambel geladen.
% Wichtig: Diese Fonts werden NICHT global auf das Dokument angewendet,
% sondern nur ueber die FDI-Font-Makros verwendet.
% Renderer=Basic ist bei alten Tengwar-/Symbolfonts oft wichtig.
%
% Der Pfad muss hier ausgeschrieben werden, weil \newfontfamily den
% Pfad schon beim Laden der Praeambel aufloest und \FDILanguage zu
% diesem Zeitpunkt noch die Sprache des Artikels ist.

% Normale FDI-Sprechblasen / Standard-Overlay
\newfontfamily{\FDIElbArialFont}{tngani_elb.ttf}[
  Path = language_files/elb/fonts_elb/,
  Renderer = Basic,
  Scale = 0.78
]

% Fette FDI-Sprechblasen / Hervorhebungen
\newfontfamily{\FDIElbArialBOLDFont}{tngani_elb.ttf}[
  Path = language_files/elb/fonts_elb/,
  Renderer = Basic,
  Scale = 0.78,
  FakeBold = 1.8
]

% FDI-Overlay-Schrift, etwas groesser fuer Comic-Texte
\newfontfamily{\FDIElbOverlayFont}{tngani_elb.ttf}[
  Path = language_files/elb/fonts_elb/,
  Renderer = Basic,
  Scale = 0.82
]

% FDI-Karolingische-Minuskel-Ersatzschrift
\newfontfamily{\FDIElbKarolingischeMinuskel}{tngani_elb.ttf}[
  Path = language_files/elb/fonts_elb/,
  Renderer = Basic,
  Scale = 0.86,
  FakeBold = 1.3
]

% Zuweisung der FDI-Font-Makros fuer diese Sprache.
% Wird von \FDISetupLanguageFonts ausgefuehrt, sobald elb gesetzt ist.
\FDISetLanguageFonts{elb}{%
  \let\KarolingischeMinuskel\FDIElbKarolingischeMinuskel
  \let\FDIOverlayFont\FDIElbOverlayFont
  \let\ArialBOLDFont\FDIElbArialBOLDFont
  \let\ArialFont\FDIElbArialFont
}
